%Texlive-full Version 3.141592-1.40.3 (Web2C 7.5.6)
%TexStudio Version 2.7.0

\documentclass[answers]{exam}
%\documentclass{exam}% choose this one to have questions only

%\newcommand*{\ANCOAMOTIVATION}{}% Uncomment to have the motivation session
\newcommand*{\ANCOACPP}{}% Uncomment to have the C++ session
\newcommand*{\ANCOAMATH}{}% Uncomment to have the math session
\newcommand*{\ANCOASTAT}{}% Uncomment to have the stat session
\newcommand*{\ANCOAFINANCE}{}% Uncomment to have the finance session
\newcommand*{\ANCOAALGO}{}% Uncomment to have the algo session

\usepackage[utf8]{inputenc}
\usepackage{graphicx}
\usepackage{subcaption}

\usepackage{lmodern}
\usepackage[a4paper]{geometry}

\usepackage{hyperref}
\usepackage{enumitem}
% \usepackage{nomencl}\makenomenclature

\usepackage{amsmath}
\usepackage{amssymb}
\usepackage{amsthm}

\usepackage{pstricks}
\usepackage{pst-node}
\usepackage{color}
\usepackage[ddmmyyyy]{datetime} 

\newcommand{\E}{\mathbb{E}}
\newcommand{\Proba}{\mathbb{P}}
\newcommand{\F}{\mathcal{F}}
\newcommand{\Bt}{B_t}
\newcommand{\DtT}{D(t,T)}
\newcommand{\RtT}{R(t,T)}
\newcommand{\PtT}{P(t,T)}
\newcommand{\PtS}{P(t,S)}
\newcommand{\LtT}{l(t,T)}
\newcommand{\YtT}{Y(t,T)}
\newcommand{\ftT}{f(t,T)}
\newcommand{\tautT}{\tau(t,T)}
\newcommand{\inttT}{\int_t^T}
\newcommand{\intt}{\int_0^t}
\newcommand{\FuncExp}{\text{exp}}
\newcommand{\FuncLn}{\text{ln}}
\newcommand{\Discount}[2]{e^{-\int_{#1}^{#2}r_u du} }
\newcommand{\bl}[1]{{\scshape  #1}}

\newcommand{\Ti}{T_{i}}
\newcommand{\Tj}{T_{j}}
\newcommand{\Ta}{T_{\alpha}}
\newcommand{\Tb}{T_{\beta}}
\newcommand{\Tii}{T_{i+1}}
\newcommand{\Pti}{P(t,T_{i})}
\newcommand{\Ptii}{P(t,T_{i+1})}
\newcommand{\Lti}{L(t,\Ti,\Tii)}
\newcommand{\Lit}{L_{i}(t)}
\newcommand{\muit}{\mu_i(t)}
\newcommand{\sigmait}{\sigma_i(t)}
\newcommand{\Wit}{W^{\Tii}(t)}
\newcommand{\Wnt}{W^{T_{N+1}}(t)}
\newcommand{\Ltj}{L(t,\Tj,\Tj)}
\newcommand{\Ljt}{L_{j}(t)}
\newcommand{\sigmajt}{\sigma_j(t)}
\newcommand{\ZCi}{ZC_{i}}
\newcommand{\ZCii}{ZC_{i+1}}
\newcommand{\Numi}{Num_{i}}
\newcommand{\Numii}{Num_{i+1}}

\newcommand{\Sab}{S_{\alpha,\beta}}
\newcommand{\Cab}{C_{\alpha,\beta}}
\newcommand{\Wab}{W_{\alpha,\beta}}
\newcommand{\sigmaab}{\sigma_{\alpha,\beta}}
\newcommand{\rhoij}{\rho_{ij}}
\newcommand{\Eab}{\mathbb{E}^{\alpha,\beta}}

\newcommand{\todo}[1]{\begin{center}\color{red} TODO : #1\end{center}}
\newcommand{\mathvec}[1]{\textbf{#1}}

\usepackage{listings}
%\lstset{language=C++,basicstyle=\scriptsize \color{green},identifierstyle=\color{orange},keywordstyle=[1]\color{blue},columns=fullflexible}
\lstset{language=C++,basicstyle=\scriptsize, keywordstyle=\color{blue}, stringstyle=\color{black},	commentstyle=\color{black} }


\begin{document}
\begin{titlepage}
%%%%%%%%%%%%%%%%%%%%%%%%%%  LOGO  %%%%%%%%%%%%%%%%%%%%%%%%%%%%%%%%%%%%%%%%%
\begin{center}
\begin{pspicture}(-7,0)(7,2)
%\psline(-7,0)(7,2)
\rput(-3.5,3){\href{http://ancoa.com/}{\includegraphics[scale=0.3]{Ancoa_Logo}}}
\end{pspicture}
\end{center}

%%%%%%%%%%%%%%%%%% DOCUMENT TITLE %%%%%%%%%%%%%%%%%%%%%%%%%
\vspace{1cm}

\begin{center}
\resizebox{14cm}{0.8cm}{Interview questions}
\end{center}
%%%%%%%%%%%%%%%%%% DOCUMENT TITLE %%%%%%%%%%%%%%%%%%%%%%%%%
\begin{center}
Version : \today
\end{center}
\vspace{2cm}
\begin{center}
The aim of the list of questions is not to require the candidate to know everything and give right answers to every questions. Those questions will be picked a few to an interview, will help us to see what area the candidate is comfortable in. We might also give all the list of question and ask the candidate to choose the one he/she is comfortable the most
\end{center}

%%%%%%%%%%%%%%%%%% ABSTRACT %%%%%%%%%%%%%%%%%%%%%%%%%
\vfill
\begin{center}
Ancoa Software Ltd \copyright \hbox{ }2012-2016.
\end{center}
\end{titlepage}
% % % % % % % % % % % % % % % % % % % % % % % % % % % % % % % % % % % % % % %
\newpage
%%%%%%%%%%%%%%%%%%%%%%%%%%%%%%%%%%%%%%%%%%%%%%%%%%%%%%%%%%%%%%%%%%%%%%%%%%%%%%%%%%%%%%%%%%
%%%%%%%%%%%%%%%%%%%%%%%%%%%%%%%%%%%%%%%%%%%%%%%%%%%%%%%%%%%%%%%%%%%%%%%%%%%%%%%%%%%%%%%%%%
\ifdefined\ANCOAMOTIVATION
\section{Motivation questions}
\renewcommand{\questionlabel}{Phone.~\thequestion.}
Since we don't have too much time for the phone interview, the main purpose is 
\begin{itemize}
  \item Check if everything mentioned in the CV is true, and the candidate is good at what he has done.
  \item Check the motivation of the candidate.
  \item Check the suspicious (discontinuous) in the career path.
\end{itemize}
Below are a few suggestions without answer :

\begin{questions}
%%%%%%%%%%%%%%%%%%%%%%%%%%%%%%%%%%%%%%%%%%
\question We're interested in your experiences/background. How did you imagine the job in Ancoa, and how do you imagine your background would contribute to Ancoa?
\begin{solution}[.2in]
	This question check not only if the candidate interest to the company/business, but also get know the creativity of the candidate. After their answer, we can also talk about how we're interested and how we imagine the candidate can add value to Ancoa.
\end{solution}

%%%%%%%%%%%%%%%%%%%%%%%%%%%%%%%%%%%%%%%%%%
\question Point randomly to a few keywords/technologies in the last experiences and ask the candidate to talk about it.

We can also pick some quick question that does not require a draft, that the candidate know or don't know. But remember to note it down, in case we invite the candidate to the face-to-face interview, we don't ask the same question.
\end{questions}
\fi

%%%%%%%%%%%%%%%%%%%%%%%%%%%%%%%%%%%%%%%%%%%%%%%%%%%%%%%%%%%%%%%%%%%%%%%%%%%%%%%%%%%%%%%%%%
%%%%%%%%%%%%%%%%%%%%%%%%%%%%%%%%%%%%%%%%%%%%%%%%%%%%%%%%%%%%%%%%%%%%%%%%%%%%%%%%%%%%%%%%%%
\ifdefined\ANCOACPP
\section{C++ programming}
\renewcommand{\questionlabel}{C++.~\thequestion]}
\begin{questions}

%%%%%%%%%%%%%%%%%%%%%%%%%%%%%%%%%%%%%%%%%%
\question How many bits in a byte?
\begin{solution}[.2in]
	8
\end{solution}

%%%%%%%%%%%%%%%%%%%%%%%%%%%%%%%%%%%%%%%%%%
\question What is inheritance?
\begin{solution}[.2in]
	Inheritance is a process of creating new classes from existing classes.	New classes inherit some of the properties and behavior of the existing classes. An existing class that is "parent" of a new class is called a base class. New class that inherits properties of the base class is called a derived class
\end{solution}

%%%%%%%%%%%%%%%%%%%%%%%%%%%%%%%%%%%%%%%%%%
\question When would you use a virtual destructor?
\begin{solution}[.2in]
Virtual destructors are useful when you can delete an instance of a derived class through a pointer to base class. Always make base classes' destructors virtual when they're meant to be manipulated polymorphically. If this is not done deleting the pointer to the base class results in undefined behaviour (most likely ‘slicing’).
\end{solution}

%%%%%%%%%%%%%%%%%%%%%%%%%%%%%%%%%%%%%%%%%%
\question What does the keyword 'override' do?
\begin{solution}[.2in]
The override keyword serves two purposes:
\begin{itemize}
	\item It shows the reader of the code that "this is a virtual method, that is overriding a virtual method of the base class."
	\item The compiler also knows that it's an override, so it can "check" that you are not altering/adding new methods that you think are overrides.
\end{itemize}
\end{solution}

%%%%%%%%%%%%%%%%%%%%%%%%%%%%%%%%%%%%%%%%%%
\question Which functions are defined for you if you do not define them in a class?
\begin{solution}[.2in]
In C++03 there are 4:
\begin{itemize}
	\item Default constructor: Declared only if no user-defined constructor is declared. Defined when used
	\item Copy constructor - declared only if the user hasn't declared one. Defined if used	
	\item Copy-assignment operator same as above	
	\item Destructor same as above	
\end{itemize}	
In C++11 there are 2 more :
\begin{itemize}
	\item Move constructor
	\item Move-assignment operator	
\end{itemize}	
\end{solution}

%%%%%%%%%%%%%%%%%%%%%%%%%%%%%%%%%%%%%%%%%%
\question What is the difference between operator new and malloc?
\begin{solution}[.2in]
\begin{center}
	\begin{tabular}{c|ll}
		\hline
		Feature & new, delete & malloc, free \\
 		\hline
		Memory allocated from & Free Store & Heap\\
 		\hline
        Returns & Fully typed pointer & void*\\
		\hline
		On failure & Throws (never erturn NULL) & Returns NULL\\
		\hline
		Required size & calculated by compiler & Must be specified in bytes\\
		\hline
		Handling arrays & Has an explicit version & Requires manual calculations\\
		\hline
		Reallocating & Not handled intuitively & Simple (no cpy constructor)\\
		\hline
		Call of reverse & Implementation defined & No\\
		\hline
		Low memory cases & Can add new memory allocator & Not handled by user code\\
		\hline
		Overridable & Yes & No\\
		\hline
		Use of con/destructor & Yes & No\\
		\hline
	\end{tabular} 
\end{center}
\end{solution}

%%%%%%%%%%%%%%%%%%%%%%%%%%%%%%%%%%%%%%%%%%
\question Name a container that has random access iterators in C++11 STL?
\begin{solution}[.2in]
	vector, array
\end{solution}

%%%%%%%%%%%%%%%%%%%%%%%%%%%%%%%%%%%%%%%%%%
\question What is not of not of the number 10. (\lstinline{!!10}).
\begin{solution}[.2in]
	\lstinline{true}
\end{solution}

%%%%%%%%%%%%%%%%%%%%%%%%%%%%%%%%%%%%%%%%%%
\question What is the difference between overriding and overloading a function?
\begin{solution}[.2in]
Overloading: changing the original types for/number of arguments in the signature of a method. Overriding: changing the original definition of a method
\end{solution}

%%%%%%%%%%%%%%%%%%%%%%%%%%%%%%%%%%%%%%%%%%
\question Which STL function can be used to sort a linked list?
\begin{solution}[.2in]
	\lstinline{std::list::sort}
\end{solution}

%%%%%%%%%%%%%%%%%%%%%%%%%%%%%%%%%%%%%%%%%%
\question What is the difference between the copy constructor and the move constructor?
\begin{solution}[.2in]
A copy constructor takes its argument by lvalue reference and a move constructor by rvalue reference.
\end{solution}

%%%%%%%%%%%%%%%%%%%%%%%%%%%%%%%%%%%%%%%%%%
\question Which keyword do you use if you want to update a member variable in a const class?
\begin{solution}[.2in]
	\lstinline{mutable}
\end{solution}

%%%%%%%%%%%%%%%%%%%%%%%%%%%%%%%%%%%%%%%%%%
\question Can you given 2 examples of design patterns and when you might use them?
\begin{solution}[.2in]
Below are a few creational design patterns :
\begin{description}
	\item[Singleton] Ensure that only one instance of a class is created and Provide a global access point to the object.
	\item[Factory] Creates objects without exposing the instantiation logic to the client and Refers to the newly created object through a common interface.
	\item[Factory Method] Defines an interface for creating objects, but let subclasses to decide which class to instantiate and Refers to the newly created object through a common interface.
	\item[Abstract Factory] Offers the interface for creating a family of related objects, without explicitly specifying their classes.
	\item[Builder] Defines an instance for creating an object but letting subclasses decide which class to instantiate and Allows a finer control over the construction process.
	\item[Prototype] Specify the kinds of objects to create using a prototypical instance, and create new objects by copying this prototype.
	\item[Object Pool] Reuses and shares objects that are expensive to create.
\end{description}
\end{solution}

%%%%%%%%%%%%%%%%%%%%%%%%%%%%%%%%%%%%%%%%%%
\question What is the size of an empty class of lambda?
\begin{solution}[.2in]
Empty class of lambda is a stateless object that take the minimum size 1.
\end{solution}

%%%%%%%%%%%%%%%%%%%%%%%%%%%%%%%%%%%%%%%%%%
\question What is the difference of vector, list and map. What is the appropriated use for each containter.
\begin{solution}[.2in]
	Candidate is expected to know vector is for use of dynamic container. List is used when the container need to delete and insert, map is used when search is needed.
\end{solution}

%%%%%%%%%%%%%%%%%%%%%%%%%%%%%%%%%%%%%%%%%%
\question What is the difference between \lstinline{class} and \lstinline{struct}.
\begin{solution}[.2in]
	Member of a class defined with the keyword \lstinline{class} are private by default. Members of a class defined with the keywords \lstinline{struct} or union are public by default.
	The keyword \lstinline{class} can be used to declare template parameters, while the \lstinline{struct} keyword cannot be so used
\end{solution}

%%%%%%%%%%%%%%%%%%%%%%%%%%%%%%%%%%%%%%%%%%
\question What is the difference of \lstinline{public}, \lstinline{private} and \lstinline{protected}.
\begin{solution}[.2in]
	Candidate is expected \lstinline{public} allow access from all other classes, \lstinline{protected} allow access from child classes, \lstinline{private} allow access only from declared \lstinline{friend} classes or functions.
\end{solution}

%%%%%%%%%%%%%%%%%%%%%%%%%%%%%%%%%%%%%%%%%%
\question What is the difference between pointers and references.
\begin{solution}[.2in]
    Candidate is expected to know references is considered as the object itself, cannot be \lstinline{null} like pointers.
\end{solution}

%%%%%%%%%%%%%%%%%%%%%%%%%%%%%%%%%%%%%%%%%%
\question What is polymorphism? Given three geometry types : Rectangle, Triangle, Circle that all need to calculate their area. How would you write those classes?
\begin{solution}[.2in]
\begin{itemize}
	\item Polymorphism is the ability to call different methods at run-time based on an object’s dynamic type. 
	\item Candidate is expected to write a BaseGeometry class, and all concrete geometries inherite from this class with a method \lstinline{virtual double calculateArea() const;}.
\end{itemize}
\end{solution}

%%%%%%%%%%%%%%%%%%%%%%%%%%%%%%%%%%%%%%%%%%
\question What is wrong with the following code?
\begin{lstlisting}
class Analytic
{
  public:
    double getSum() const
    {
      _bufferSum=0;
      for(auto const& elem : mTimeSerie)
        _bufferSum+=elem;
      return _bufferSum;
    }

private:
    std::vector<double> mTimeSerie;
    double _bufferSum;
};

\end{lstlisting}
\begin{solution}[.2in]
\lstinline{const} member methods cannot modify member variables. Candidate is expected to suggest to set \lstinline{_bufferSum} to \lstinline{mutable}, or use local variable inside the for loop instead of member variable.
\end{solution}

%%%%%%%%%%%%%%%%%%%%%%%%%%%%%%%%%%%%%%%%%%
\question What is wrong with the following code?
\begin{lstlisting}
void printNode(Node* pNode)
{
    std::cout << pNode->getValue() << std::endl;
}
\end{lstlisting}
\begin{solution}[.2in]
Possible null pointer dereference
\end{solution}

%%%%%%%%%%%%%%%%%%%%%%%%%%%%%%%%%%%%%%%%%%
\question What is wrong with the following code?
\begin{lstlisting}
void SomeClass::setProperties(const std::vector<Property>& rPropertiesVec)
{
    if(rPropertiesVec.empty())
      return;

    mFirstProperty = rPropertiesVec.at(0);
    mSecondProperty = rPropertiesVec.at(1);
}
\end{lstlisting}
\begin{solution}[.2in]
Possible out of bounds container access
\end{solution}

%%%%%%%%%%%%%%%%%%%%%%%%%%%%%%%%%%%%%%%%%%
\question Write a function in C++ to swap the values of two variables of the same type.
\begin{solution}[.2in]
\begin{lstlisting}
template <class T>
void swap(T& rArg1, T& rArg2)
{
    T temp = rArg1;
    rArg1 = rArg2;
    rArg2 = temp;
}
\end{lstlisting}
\end{solution}

%%%%%%%%%%%%%%%%%%%%%%%%%%%%%%%%%%%%%%%%%%
\question What is wrong with the following code?
\begin{lstlisting}
double calculateMagicSum(const std::vector<double>& rNumbersVec)
{
    double sum = 0.0;
    for(int i = 0; i < rNumbersVec.size(); ++i)
    {
      MagicCalc* pCalc = new MagicCalc;
      const double magicValue = pCalc->calculate(rNumbersVec.at(i));
      sum += magicValue;
    }
    return sum;
}
\end{lstlisting}
\begin{solution}[.2in]
The MagicCalc objects are leaked
\end{solution}

%%%%%%%%%%%%%%%%%%%%%%%%%%%%%%%%%%%%%%%%%%
\question What is wrong with the following code?
\begin{lstlisting}
class Base
{
  public:
    Base();

  private:
    virtual void someFunction();

    int mBaseVar;
};

class Derived : public Base
{
  public:
    Derived();

  private:
    virtual void someFunction();

    int mDerivedVar;
};
\end{lstlisting}
\begin{solution}[.2in]
virtual destructor is not declared
\end{solution}

%%%%%%%%%%%%%%%%%%%%%%%%%%%%%%%%%%%%%%%%%%
\question Write the body of the function 'printTree' to traverse a tree and print each node’s value.
\begin{lstlisting}
struct TreeNode
{
    void print() const;

    void* mpData;
    std::vector<TreeNode*> mChildren;
};

void printTree(TreeNode* pRootNode) const
{
    ...
}
\end{lstlisting}
\begin{solution}[.2in]
\begin{lstlisting}
void printTree(TreeNode* pRootNode) const
{
    if(!pRootNode)
      return;

    // Print this node
    pRootNode->print();

    // Print the child nodes
    std::vector<TreeNode*>::const_iterator citer;
    for(citer = children.begin();citer != children.end();++citer)
    {
      printTree(*citer);
    }
}
\end{lstlisting}
\end{solution}

%%%%%%%%%%%%%%%%%%%%%%%%%%%%%%%%%%%%%%%%%%
\question Suppose you have an array of integers, in which one or more integers are duplicated.
Write the body of the following function to find all the duplicates.
\begin{lstlisting}
void findDuplicatesIndices(const std::vector<int>& rIntsVec, 
std::vector<int>& rDuplicatesIndicesVec) const
{
    //Fill here
}
\end{lstlisting}
\begin{solution}[.2in]
\begin{lstlisting}
void findDuplicatesIndices(const std::vector<int>& rIntsVec,
std::vector<int>& rDuplicatesIndicesVec) const
{
    rDuplicatesIndicesVec.clear();

    std::set<int> seenSoFar;
    for(int i = 0; i < rIntsVec.size(); ++i)
    {
      int val = rIntsVec[i];
      if(seenSoFar.find(val) != seenSoFar.end())
      rDuplicatesIndicesVec.push_back(i);
      seenSoFar.insert(val);
    }
}
\end{lstlisting}
\end{solution}

%%%%%%%%%%%%%%%%%%%%%%%%%%%%%%%%%%%%%%%%%%
\question What is wrong with the following code?
\begin{lstlisting}
void removeDuplicates(std::vector<int>& rIntVec)
{
    std::set<int> values;
    std::vector<int>::iterator iter;
    for(iter = std::begin(rIntVec); iter != std::end(rIntVec); ++iter)
    {
      if(values.count(*iter))
        rIntVec.erase(iter);
      else
        values.insert(*iter);
    }
}
\end{lstlisting}
\begin{solution}[.2in]
Function which deletes objects via an iterator which invalidates the iterator. How should this code be fixed?
\begin{lstlisting}
void removeDuplicates(std::vector<int>& rIntVec)
{
    std::set<int> values;
    std::vector<int>::iterator iter;
    for(iter = std::begin(rIntVec); iter != std::end(rIntVec); ++iter)
    {
      if(values.count(*iter))
        iter = rIntVec.erase(iter) - 1;
      else
        values.insert(*iter);
}
}
\end{lstlisting}
\end{solution}

%%%%%%%%%%%%%%%%%%%%%%%%%%%%%%%%%%%%%%%%%%
\question Write the body of the function 'reverse' to reverse a singly linked list.
\begin{lstlisting}
struct ListNode
{
    ListNode* mpNext; // null indicates end of list
    void* mpData;
};

ListNode* reverse(ListNode* pHeadNode)
{
    // Reverse the list defined by pHeadNode
    // Return the head node of the reversed list
    ...
}
\end{lstlisting}
\begin{solution}[.2in]
\begin{lstlisting}
ListNode* reverse(ListNode* pHeadNode)
{
    if(!pHeadNode)
      return 0;

    ListNode* pCurrentNode = pHeadNode;
    ListNode* pPrevNode = 0;

    while(pCurrentNode)
    {
       // Swap the 'next' pointer
       ListNode* pNextNode = pCurrentNode->mpNext;
       pCurrentNode->mpNext = pPrevNode;

       // Move on to the next node
       pPrevNode = pCurrentNode;
       pCurrentNode = pNextNode;
    }
    return pPrevNode;
}
\end{lstlisting}
\end{solution}

%%%%%%%%%%%%%%%%%%%%%%%%%%%%%%%%%%%%%%%%%%
\question What would be printed to the console when this code is run?
\begin{lstlisting}
int main()
{
    int a, b, c = 3;
    a = b = 1;
    ++b += c++ + (a+=3);
    b *= 3;
    while(--c) a - 2;
    a += c *= 2, ++b, b - 4;
    (c ? a : b) += b ? !!a : !a;
    std::cout << a << "," << b << "," << c << std::endl;
}
\end{lstlisting}
\begin{solution}[.2in]
4, 29, 0
\begin{lstlisting}
int main()
{
    int a, b, c = 3; // After exec: a = ?, b = ?, c = 3.
    a = b = 1; //  After exec: a = 1, b = 1, c = 3.
    ++b += c++ + (a+=3); // After exec: a = 4, b = 9, c = 4.
    b *= 3; // After exec: a = 4, b = 27, c = 4.
    while(--c) a - 2; // After exec: a = 4, b = 27, c = 0.
    a += c *= 2, ++b, b - 4; //  After exec: a = 4, b = 28, c = 0.
    (c ? a : b) += b ? !!a : !a; // After exec: a = 4, b = 29, c = 0.
    std::cout << a << "," << b << "," << c << std::endl;
}
\end{lstlisting}
\end{solution}

\end{questions}
\fi
%%%%%%%%%%%%%%%%%%%%%%%%%%%%%%%%%%%%%%%%%%%%%%%%%%%%%%%%%%%%%%%%%%%%%%%%%%%%%%%%%%%%%%%%%%
%%%%%%%%%%%%%%%%%%%%%%%%%%%%%%%%%%%%%%%%%%%%%%%%%%%%%%%%%%%%%%%%%%%%%%%%%%%%%%%%%%%%%%%%%%
\ifdefined\ANCOAMATH
\newpage
\section{Mathematics}
\renewcommand{\questionlabel}{Math.~\thequestion ]}
\begin{questions}
%%%%%%%%%%%%%%%%%%%%%%%%%%%%%%%%%%%%%%%%%%
\question How to calculate the sum from 1 to 20.
\begin{solution}[.2in]
	The sum from 1 to $N$ is $\frac{1}{2}N(N+1)$. The sum from 1 to 20 is 210.
\end{solution}

%%%%%%%%%%%%%%%%%%%%%%%%%%%%%%%%%%%%%%%%%%
\question In a debugging situation : Customer give their data as a time series of prices, and run our analytic tool to calculate the average of prices. The client then look at the trading information on internet and complain the average should be around 270 while the analytic give 243.05 as result. How do you would you answer to the customer? 
\[
\left[ 
\begin{array}{c}
270.33 \\
269.87 \\
270.17 \\
0.0027 \\
270.19 \\
269.97 \\
270.33 \\
270.33 \\
269.42 \\
269.92
\end{array}
\right]
\hspace{2cm}
\text{analytics result 243.05327}
\]
\begin{solution}[.2in]
	The exact sum of the vector is 243.05327. The candidate is expected to see first there're nearly a zero in the vector. And know to approximatively calculate the average by $~ \frac{270*9}{10}$
\end{solution}

%%%%%%%%%%%%%%%%%%%%%%%%%%%%%%%%%%%%%%%%%%
\question How to evaluate matrix vector multiplication? How to evaluate bilinear on 2 vectors (example 2 dimension case).
\[
\left[ 
\begin{array}{cc}
a_{11} & a_{12}\\
a_{21} & a_{22}
\end{array} 
\right]
\left[ 
\begin{array}{c}
x_{1}\\
x_{2}
\end{array} 
\right]
=?
\hspace{1cm}
\left[ 
\begin{array}{cc}
x_{1} & x_{2}
\end{array} 
\right]
\left[ 
\begin{array}{cc}
a_{11} & a_{12}\\
a_{21} & a_{22}
\end{array} 
\right]
\left[ 
\begin{array}{c}
x_{1}\\
x_{2}
\end{array} 
\right]
=?
\]
\begin{solution}[.2in]
\[
\left[ 
\begin{array}{cc}
a_{11} & a_{12}\\
a_{21} & a_{22}
\end{array} 
\right]
\left[ 
\begin{array}{c}
x_{1}\\
x_{2}
\end{array} 
\right]
=
\left[ 
\begin{array}{c}
a_{11}x_1 + a_{12}x_2\\
a_{21}x_1 + a_{22}x_2
\end{array} 
\right]
\]
And
\[
\left[ 
\begin{array}{cc}
x_{1} & x_{2}
\end{array} 
\right]
\left[ 
\begin{array}{cc}
a_{11} & a_{12}\\
a_{21} & a_{22}
\end{array} 
\right]
\left[ 
\begin{array}{c}
x_{1}\\
x_{2}
\end{array} 
\right]
=
a_{11}x^2_1 + (a_{12}+a_{21})x_1x_2 + a_{22}x^2_2
\]
\end{solution}

%%%%%%%%%%%%%%%%%%%%%%%%%%%%%%%%%%%%%%%%%%
\question Providing a rectangle in 2D plan with coordinates of 4 corners. How to sample 1 millions points uniformly inside the rectangle.
\begin{solution}[.2in]
Todo : Candidate is expected to know to distribute 4 points uniformly inside the rectangle. Then generalize the idea to 1M points. The most common way is to be able to calculate the transformation matrix, transform an arbitrary rectangle to the reference rectangle (unit square). The invert of this matrix transform the unit square to the rectangle. Uniformly distribute points inside the unit square is easy, then transform them by mean of transformation matrix.
\end{solution}

%%%%%%%%%%%%%%%%%%%%%%%%%%%%%%%%%%%%%%%%%%
\question What is a positive definite symmetric matrix.
\begin{solution}[.2in]
A matrix $A$ is positive definite iff 
\[
\mathvec{x}^t A \mathvec{x} = 0 \Longleftrightarrow \mathvec{x}=0
\]
\end{solution}

%%%%%%%%%%%%%%%%%%%%%%%%%%%%%%%%%%%%%%%%%%
\question How to calculate the norm of a $N$-dimensional vector. Do you know other norm? Which condition should satisfied for a norm.
\begin{solution}[.2in]
Given $\mathvec{x}=[x_1,x_1 ... x_n] \in \mathbb{R}^n$
\begin{itemize}
	\item Norm $L^2$ (Euclidian norm) is : $||\textbf{x}||_{2} = \sqrt{\sum x^2_i} $
	\item Norm $L^1$ is : $||\textbf{x}||_{1} = \sum |x_i|$
	\item Norm $L^\infty$ is : $||\textbf{x}||_{\infty} = \textbf{max} |x_i|$
\end{itemize}
\end{solution}

%%%%%%%%%%%%%%%%%%%%%%%%%%%%%%%%%%%%%%%%%%
\question What is a eigenvalues and eigenvectors. Method to find eigenvalues. Is that easy to find eigenvectors?
\begin{solution}[.2in]
Eigenvalues are associated with eigenvectors, describe spectrum spaces of a squared matrix $\mathvec{A}$ by equations 
\[
\lambda\mathvec{x} = \mathvec{A}\mathvec{x}
\]
$\lambda$ is the eigenvalue, $\mathvec{x}$ is the associated eigenvector. Indeed eigenvector is in the kernel of the application (matrix) $\mathvec{A}-\lambda\mathvec{I}$, so eigenvalues can be found by resolve the equation
\[
\text{det}(\mathvec{A}-\lambda\mathvec{I})=0
\]
Eigenvectors or eingenspace can be manually calculated for small matrices. For big matrices, there are a few algorithm to find eigenvectors, implemented in C++.
\end{solution}

%%%%%%%%%%%%%%%%%%%%%%%%%%%%%%%%%%%%%%%%%%
\question Why and how to invert a matrix? Example invert a 2x2 matrix 
\[
\mathvec{A}=
\left[ 
\begin{array}{cc}
a_{11} & a_{12}\\
a_{21} & a_{22}
\end{array} 
\right]
\]
Some methods to invert very big matrix 1Mx1M
\begin{solution}[.2in]
There are two methods to revert a small matrix : or calculate determinant, adjoint matrices, then apply the invert formula an simple way is to write an unknown matrix $\mathvec{X}$ and solve the equation $\mathvec{X}\mathvec{A}=\mathvec{I}_2$. The result has to be :
\[
\mathvec{A}^{-1}=
\frac{1}{a_{11}a_{22}-a_{12}a_{21}}
\left[ 
\begin{array}{cc}
a_{22} & -a_{12}\\
-a_{21} & a_{11}
\end{array} 
\right]
\]
Two basical methods to invert matrices are LU, Conjugated Gradient. The first is direct method which require to store the whole matrix in memory, the second is iterative, don't need to store the matrix in memory, but implement the operation $\mathvec{x}+=\mathvec{A}d\mathvec{x}$
\end{solution}

%%%%%%%%%%%%%%%%%%%%%%%%%%%%%%%%%%%%%%%%%%
\question How to check if two rectangles in 2D space overlaps?
\begin{solution}[.2in]
First, two rectangles are covered by 2 circles. If those two circles are not overlaped, the two rectangles are not. If the two circles are overlaped, then check if a corner of one rectangle is inside other rectangle. To check if a point $M$ is inside the rectangle $ABCD$, check the projection of $M$ to $AB$ and $AD$, which is 
\[
\frac{AM \cdot AB}{AB \cdot AB} \leq 1  
\cap 
\frac{AM \cdot AD}{AD \cdot AD} \leq 1  
\]
\end{solution}

%%%%%%%%%%%%%%%%%%%%%%%%%%%%%%%%%%%%%%%%%%
\question Intuition of how google map work, give the idea of how to visualize a matrix dimension 1Bx1B?
\begin{solution}[.2in]
The idea could be building a hierarchical bounding volume. At root level, matrix is a single rectangle, with a single color. At second level, the matrix can be split in two, and the half having more non-null element will have a darker color. At third level, the two half matrix are split again, each into two part.... Depending of the zoom of the viewer, one can calculate the detail level and render the image.
\end{solution}

%%%%%%%%%%%%%%%%%%%%%%%%%%%%%%%%%%%%%%%%%%
\question Simple way to parallelise Matrix * Vector
\begin{solution}[.2in]
Expected to know multiplication by blocks.
\end{solution}

%%%%%%%%%%%%%%%%%%%%%%%%%%%%%%%%%%%%%%%%%%
\question What are two orthogonal vectors? Calculate orthogonal projection of $\mathvec{x}$ into $\mathvec{y}$. Test with a simple example of dimension 2. 
\begin{solution}[.2in]
Two vectors are orthogonal iff $<\mathvec{x},\mathvec{y}>=0$. Projection of $\mathvec{x}$ on $\mathvec{y}$ is
\[
\mathvec{x} - \frac{<\mathvec{x},\mathvec{y}>}{||\mathvec{y}||^2} \mathvec{y}
\]
Test in 2D with $\mathvec{x}^t=[1,1]$ and $\mathvec{y}^t=[-1,1]$
\end{solution}

%%%%%%%%%%%%%%%%%%%%%%%%%%%%%%%%%%%%%%%%%%
\question Draw graphic of functions $y=2x-1$, $y=|x|$, $y=e^x$, $y=ln(x)$, $y=cos(x)$, $y=sin(x)$.
\begin{solution}[.2in]
	todo
\end{solution}

%%%%%%%%%%%%%%%%%%%%%%%%%%%%%%%%%%%%%%%%%%
\question Draw graphic of unit bull with difference distance definition L2, L1, LInf.
\begin{solution}[.2in]
	todo
\end{solution}

%%%%%%%%%%%%%%%%%%%%%%%%%%%%%%%%%%%%%%%%%%
\question What is Newton Raphson method. How to test if the Newton Raphson implementation give correct result
\begin{solution}[.2in]
Newton Raphson is a numerical method to find solution of a function $f(x)=a$. One can test the implementation by giving a point $\hat{x}$, calculate $\hat{a}=f(\hat{x})$, and resolve the equation $f(x)=\hat{a}$. If the found solution is close, very closed to $\hat{x}$, the algorithm is well implemented.
\end{solution}

%%%%%%%%%%%%%%%%%%%%%%%%%%%%%%%%%%%%%%%%%%
\question What is linear interpolation? Given 2 points, how to calculate the interpolated point? We might need higher degree polynomial to interpolate, why? Given polynomial of n degree, what is the minimum number of points are needed to do a correct interpolation?
\begin{solution}[.2in]
	todo
\end{solution}

%%%%%%%%%%%%%%%%%%%%%%%%%%%%%%%%%%%%%%%%%%
\question Given a vector of bool : b1,b2… bn. Write a function to check OR of all, AND of all, negation OR of all, negation AND of all elements.
\begin{solution}[.2in]
Iterate through the vector, 
\begin{itemize}
	\item the first \lstinline{true} encountered will return \lstinline{true} for the OR function
	\item the first \lstinline{false} encountered will return \lstinline{false} for the AND function
\end{itemize}
\end{solution}


\end{questions}
\fi
%%%%%%%%%%%%%%%%%%%%%%%%%%%%%%%%%%%%%%%%%%%%%%%%%%%%%%%%%%%%%%%%%%%%%%%%%%%%%%%%%%%%%%%%%%
%%%%%%%%%%%%%%%%%%%%%%%%%%%%%%%%%%%%%%%%%%%%%%%%%%%%%%%%%%%%%%%%%%%%%%%%%%%%%%%%%%%%%%%%%%
\ifdefined\ANCOASTAT
\newpage
\section{Statistics}
\renewcommand{\questionlabel}{Stat.~\thequestion ]}
\begin{questions}

%%%%%%%%%%%%%%%%%%%%%%%%%%%%%%%%%%%%%%%%%%
\question How to calculate mean, variance, and correlation ($X,Y$).
\begin{solution}[.2in]
\begin{description}
	\item[sample mean] $\bar{x} = \frac{1}{N}\sum x_i$
	\item[sample variance] $\bar{\sigma}_{x}^2 = \frac{1}{N-1} \sum(x_i - \bar{x})^2$
	\item[sample covariance] $\bar{\sigma}_{xy}^2 = \frac{1}{N-1} \sum(x_i - \bar{x})(y_j - \bar{y})$
\end{description}
\end{solution}

%%%%%%%%%%%%%%%%%%%%%%%%%%%%%%%%%%%%%%%%%%
\question What does it mean two variables are highly correlated. Is uncorrelated say independant? Example of variable uncorrelated but dependant
\begin{solution}[.2in]
	$y=|x|$
\end{solution}

%%%%%%%%%%%%%%%%%%%%%%%%%%%%%%%%%%%%%%%%%%
\question Definition of the covariance matrix. How to get the correlation matrix from covariance matrix? What are properties of covariance matrix and correlation matrix.
\begin{solution}[.2in]
\begin{description}
	\item[Covariance Matrix] Given a set of random variables $X_i$, covariance matrix is $M_ij$ equal to the covariance of $X_i$ and $X_j$
	\item[Correlation Matrix] Correlation matrix can get by Cholesky factorization from covariance matrix. 
	\item[Properties] They're both symetric semi positive definite. Correlation matrix has values bounded by $[-1,1]$
\end{description}
\end{solution}

%%%%%%%%%%%%%%%%%%%%%%%%%%%%%%%%%%%%%%%%%%
\question What is the difference of mean vs median.
\begin{solution}[.2in]
In an sorted array of real values, mean is the average of all elements, while median is the value of the element in middle of the array. In a skewed array, mean is different than median.
\end{solution}

%%%%%%%%%%%%%%%%%%%%%%%%%%%%%%%%%%%%%%%%%%
\question What is the density of a variable? discrete case, continuous case. What condition a function should satisfy a density?
\begin{solution}[.2in]
A density function should be positive, integrable and it's integral (sum) equal to 1.
\end{solution}

%%%%%%%%%%%%%%%%%%%%%%%%%%%%%%%%%%%%%%%%%%
\question What is the cumulative distribution function? What properties it has?
\begin{solution}[.2in]
Cumulative distribution is the function describing probability of a rv $f(t) = P(X\leq t)$. Its a increasing function, growing from zero to 1.
\end{solution}

%%%%%%%%%%%%%%%%%%%%%%%%%%%%%%%%%%%%%%%%%%
\question What is a Gaussian variable? what are their properties? Example . How the shape of the gaussian density should be?  Draw a Gaussian density shape with big variance vs small variance.
\begin{solution}[.2in]
	todo
\end{solution}

%%%%%%%%%%%%%%%%%%%%%%%%%%%%%%%%%%%%%%%%%%
\question What is the idea of regression? linear regression? example from a scatter plot, how to draw intuitively a regression? How to test a regression without having real data?
\begin{solution}[.2in]
Regression is the estimation of a rv in relation to other rv. Mathematically, that estimate the conditional expectation. Linear regression make hypothesis two variables are linearly linked and estimate this linearity. Regression is indeed a parametric estimation. One should know how to virtually test the regression with simulated data, before applying to real data.
\end{solution}

%%%%%%%%%%%%%%%%%%%%%%%%%%%%%%%%%%%%%%%%%%
\question If a simulator simulate 1M values of uniform law vs Gaussian law, what should be the layout of values between -1 and 1? How to test if the gaussian simulation is correct?
\begin{solution}[.2in]
Uniform simulation will give the set of points uniformly distributed in the segment, Gaussian simulation will give the set of points more dense in the center of the segment. To test if the Gaussian simulation is correct, one should calculate sample mean and sample variance, and compare with the Gaussian parameters.
\end{solution}
\end{questions}
\fi
%%%%%%%%%%%%%%%%%%%%%%%%%%%%%%%%%%%%%%%%%%%%%%%%%%%%%%%%%%%%%%%%%%%%%%%%%%%%%%%%%%%%%%%%%%
%%%%%%%%%%%%%%%%%%%%%%%%%%%%%%%%%%%%%%%%%%%%%%%%%%%%%%%%%%%%%%%%%%%%%%%%%%%%%%%%%%%%%%%%%%
\ifdefined\ANCOAFINANCE
\newpage
\section{Finance}
\renewcommand{\questionlabel}{Phi.~\thequestion ]}
\begin{questions}
%%%%%%%%%%%%%%%%%%%%%%%%%%%%%%%%%%%%%%%%%%
\question Knowledge of any finance instrument? stock, forward, options, swap, bond …?
\begin{solution}[.2in]
\begin{itemize}
	\item For Stock, candidate is expected to know about dividend, when dividends pay out, stock price drop
    \item For Forward, candidate is expected to know the fair price formula $F=e^{rT}S$
    \item For Option, candidate is expected to know the important of volatility (the only value unobservable).
    \item For Bond, candidate is expected to know how to price a bond given Yield curve.
\end{itemize}
\end{solution}

%%%%%%%%%%%%%%%%%%%%%%%%%%%%%%%%%%%%%%%%%%
\question What is an order book? What is a market order, limit order, stop limit order, stop loss order?
\begin{solution}[.2in]
\begin{itemize}
	\item Market order is an order matching immediately at the best price
	\item Limit order is an order matching at user specified price
	\item Stop loss order is an order when price move to a specific level, become a market order.
	\item Stop limit order is an order, when price move to a specific level, become limit order.
\end{itemize}
\end{solution}

%%%%%%%%%%%%%%%%%%%%%%%%%%%%%%%%%%%%%%%%%%
\question Knowledge about liquidity? liquidity risk?
\begin{solution}[.2in]
Liquidity indicator define how easy participant can trade. Liquidity risk is the risk when investers want to trade, but cannot trade at the desired moment, with the desired volume.
\end{solution}

%%%%%%%%%%%%%%%%%%%%%%%%%%%%%%%%%%%%%%%%%%
\question What is the volatility? How to calculate?
\begin{solution}[.2in]
Volatility is a measure of variation of price. There are historical volatility and implied volatility. Historical volatility is calculated by statistical methods, implied volatility is calculated by reverting a model's pricer.
\end{solution}

%%%%%%%%%%%%%%%%%%%%%%%%%%%%%%%%%%%%%%%%%%
\question Different of Stock market, Commodity market, Interest Rate market, FX market?
\begin{solution}[.2in]
	todo
\end{solution}

%%%%%%%%%%%%%%%%%%%%%%%%%%%%%%%%%%%%%%%%%%
\question What is a Yield curve? How to have the Yield curve?
\begin{solution}[.2in]
\begin{itemize}
	\item Yield curve is the curve describing the interest rate at the present moment in function of maturity.
    \item Yield curve is quoted by data provider, or can be calculated through boostraping method or optimization method, with a bunch of Fixed Income instruments (Bonds, Forward, Swap).
\end{itemize}
\end{solution}

%%%%%%%%%%%%%%%%%%%%%%%%%%%%%%%%%%%%%%%%%%
\question Knowledge of any kind of market manipulation? If you have the super power that you can move the price of the market, what would you do?
\begin{solution}[.2in]
Candidate is expected to say : move down the price, buy the asset, move up the price, sell the asset.
\end{solution}

%%%%%%%%%%%%%%%%%%%%%%%%%%%%%%%%%%%%%%%%%%
\question If the US interest rate raise, what will be expected the FX of pound price against dollar?
\begin{solution}[.2in]
Interest rate increases means the value of the money increase. By comparing to other money, dollar will be stronger, so pound will be cheaper.
\end{solution}

%%%%%%%%%%%%%%%%%%%%%%%%%%%%%%%%%%%%%%%%%%
\question An invester want to sell a huge amount of a certain security. He split in small chunks and put in the order books at different choice of time to optimize his executions. Why he need to do so? How the price will behave, why?
\begin{solution}[.2in]
Following naturally the demand-offer balance, if there are more sellers, the price will move down. Invester should split a huge trade into small chunk to minimize the price impact. 
\end{solution}

\end{questions}
\fi

%%%%%%%%%%%%%%%%%%%%%%%%%%%%%%%%%%%%%%%%%%%%%%%%%%%%%%%%%%%%%%%%%%%%%%%%%%%%%%%%%%%%%%%%%%
%%%%%%%%%%%%%%%%%%%%%%%%%%%%%%%%%%%%%%%%%%%%%%%%%%%%%%%%%%%%%%%%%%%%%%%%%%%%%%%%%%%%%%%%%%
\ifdefined\ANCOAALGO
\newpage
\section{Algorithms/optimal Programming}
\renewcommand{\questionlabel}{Opt.~\thequestion ]}
\begin{questions}

%%%%%%%%%%%%%%%%%%%%%%%%%%%%%%%%%%%%%%%%%%
\question What is a binary search tree, what is the general complexity in BST? What is its default
\begin{solution}[.2in]
Binary Search Tree is the data structure that sort the data on the left or right child node relevant to the ordering of data. The tree data structure help to improve the complexity to $\text{log}(N)$ instead of $O(N)$. It can be a bad situation if the data put in BST is not well balanced, it can be improved by the Red-Black Tree.
\end{solution}

%%%%%%%%%%%%%%%%%%%%%%%%%%%%%%%%%%%%%%%%%%
\question How to calculate the running average in a context of real time application? (array append new element continuously)
\begin{solution}[.2in]
	To avoid recalculate the mean of the whole array everytime the array appends new element, one should store a buffer of sum and the number of element. When the array append a new element, just update the buffer.
	\paragraph{} How about running average of a running queue? Use the same approach, do not recalculate the whole queue when the data queue push or pop, just update the sum and the number of element.
\end{solution}

%%%%%%%%%%%%%%%%%%%%%%%%%%%%%%%%%%%%%%%%%%
\question As a scientific developer, what might upset you in this code 
\begin{lstlisting}
class SumCalculator
{
public:
	double calc(std::vector<double> timeseries)
	{
		double sum(0.);
		for(size_t i=0;i<timeseries.size();i++)
		    sum+=timeseries[i];
		return sum;
	}
};
\end{lstlisting}
\begin{solution}[.2in] 
\begin{lstlisting}
class SumCalculator
{
public:
	double calc(const std::vector<double>& timeseries) const
	{
		double sum(0.);
		const size_t N = timeseries.size();
		for(size_t i=0;i<N;++i)
		   sum+=timeseries[i];
		return sum;
	}
};
\end{lstlisting}
\begin{itemize}
	\item The copy in function argument is critical
	\item The \lstinline{const} in function member
	\item Prefix increment operator in the \lstinline{for} loop
	\item Use of \lstinline{N} instead of calling \lstinline{vector.size()}
\end{itemize}
Show a picky developer.
\end{solution}

\end{questions}
\fi
%\printnomenclature % % % % % % % % % % %
%\bibliographystyle{siam}
%\bibliography{interview.bib}
\end{document}
